% !TeX program = pdflatex
\documentclass[aspectratio=169,xcolor={dvipsnames,table}]{beamer}

\input{../../_en_preambulo_beamer}
% Comandos y macros estadísticas compartidas para las presentaciones Beamer en inglés (Metropolis)

% Conjuntos numéricos básicos
\newcommand{\R}{\mathbb{R}}
\newcommand{\N}{\mathbb{N}}
\newcommand{\Q}{\mathbb{Q}}
\newcommand{\Z}{\mathbb{Z}}
\newcommand{\set}[1]{\left\{ #1 \right\}}
\newcommand{\abs}[1]{\left|#1\right|}
\newcommand{\norm}[1]{\left\| #1 \right\|}

% Comandos estadísticos y probabilísticos del curso
\newcommand{\Var}{\operatorname{Var}}
\newcommand{\cov}{\operatorname{Cov}}
\newcommand{\corr}{\rho}
\newcommand{\comb}[2]{\begin{pmatrix} #1 \\ #2 \end{pmatrix}}
\newcommand{\s}{\sigma}
\newcommand{\particion}{\mathfrak{P}}
\newcommand{\card}[1]{\#\left( #1 \right)}
\newcommand{\seq}[3]{\set{#1}_{#2}^{#3}}
\newcommand{\E}{\mathbb{E}}
\newcommand{\Prob}{\mathbb{P}}

% Entornos pedagógicos y cajas en inglés académico natural (soportando tanto nombres en inglés como en español)
\newenvironment{discussion}{%
  \begin{alertblock}{Discussion Question}%
}{%
  \end{alertblock}%
}
\newenvironment{discusion}{%
  \begin{alertblock}{Discussion Question}%
}{%
  \end{alertblock}%
}

\newenvironment{reflection}{%
  \begin{alertblock}{Classroom Reflection}%
}{%
  \end{alertblock}%
}
\newenvironment{reflexion}{%
  \begin{alertblock}{Classroom Reflection}%
}{%
  \end{alertblock}%
}

\newenvironment{paradox}{%
  \begin{alertblock}{Exploratory Paradox}%
}{%
  \end{alertblock}%
}
\newenvironment{paradoja}{%
  \begin{alertblock}{Exploratory Paradox}%
}{%
  \end{alertblock}%
}

\newenvironment{motivation}{%
  \begin{exampleblock}{Motivation: Data Science in Practice}%
}{%
  \end{exampleblock}%
}
\newenvironment{motivacion}{%
  \begin{exampleblock}{Motivation: Data Science in Practice}%
}{%
  \end{exampleblock}%
}

\newenvironment{quickchallenge}{%
  \begin{block}{Quick Team Challenge}%
}{%
  \end{block}%
}
\newenvironment{retoclase}{%
  \begin{block}{Quick Team Challenge}%
}{%
  \end{block}%
}


\title{Normal Distribution}
\subtitle{Section 04.05 --- Gaussian Bell, Z-Score Standardization, and CLT}
\author[J. Castillo Colmenares]{Juliho Castillo Colmenares}
\institute[Tec de Monterrey]{Tecnologico de Monterrey}
\date{\vspace{-1.5cm}}

\begin{document}

% SLIDE 01: Title Page
\begin{frame}[plain]
	\titlepage
\end{frame}

% SLIDE 02: Roadmap
\begin{frame}{Roadmap --- Chapter 04: Continuous Random Variables}
	\vspace{-0.15cm}
	\begin{block}{Curricular Structure of Unit 3}
		\footnotesize
		\begin{enumerate}\itemsep=0.03cm
			\item \textcolor{gray}{Section 04.01: PDF and Continuous Support}
			\item \textcolor{gray}{Section 04.02: Continuous CDF and Quantiles}
			\item \textcolor{gray}{Section 04.03: Expectation, Variance, and Continuous LOTUS}
			\item \textcolor{gray}{Section 04.04: Continuous Uniform Distribution}
			\item \textbf{\textcolor{TecRojo}{Section 04.05: Normal Distribution}} $\leftarrow$ \emph{Current focus}
			\item \textcolor{gray}{Section 04.06: Exponential Distribution and Memoryless Processes}
			\item \textcolor{gray}{Section 04.07: Gamma, Beta, and Weibull Distributions}
		\end{enumerate}
	\end{block}
	\vspace{-0.2cm}
	\begin{alertblock}{Learning Objective}
		\scriptsize
		Introduce the Normal distribution as the most important of all continuous statistics, master Z-score standardization and the 68-95-99.7 rule, apply the MGF to prove sum-closure properties, and establish the Central Limit Theorem as the theoretical justification for the universality of the Normal.
	\end{alertblock}
\end{frame}

% SLIDE 03: Motivation
\begin{frame}{The Gaussian Bell: Universality of the Normal}
	\footnotesize
	\begin{motivacion}[Why the Normal distribution?]
		The Normal distribution $N(\mu, \sigma^{2})$ is arguably the most important in all of statistics. Its bell shape models measurement errors, signal noise, human heights, exam scores, and many other natural phenomena. By the Central Limit Theorem, the sum (or average) of i.i.d. variables converges to the Normal regardless of the original distribution.
	\end{motivacion}
	\vspace{0.15cm}
	\pause
	\begin{itemize}\itemsep=0.1cm
		\item \textbf{Bell shape:} Symmetric around $\mu$ with dispersion controlled by $\sigma$. \pause
		\item \textbf{Standardization:} Any Normal reduces to $N(0,1)$ via $Z = (X-\mu)/\sigma$. \pause
		\item \textbf{Universality:} Appears in the CLT as the limit of sums; foundation of statistical inference.
	\end{itemize}
\end{frame}

% SLIDE 03B: Definition and Z-Score
\begin{frame}{Definition and Z-Score Standardization}
	\vspace{-0.15cm}
	\begin{columns}[T]
		\begin{column}{0.48\textwidth}
			\begin{block}{PDF and Symmetry}
				\footnotesize
				$f(x) = \frac{1}{\sigma\sqrt{2\pi}} \exp\left(-\frac{(x-\mu)^{2}}{2\sigma^{2}}\right)$. Bell-shaped, symmetric around $\mu$. Support $\mathbb{R}$. Mean, median, and mode all equal $\mu$.
			\end{block}
		\end{column}
		\begin{column}{0.48\textwidth}
			\begin{alertblock}{Z-Score and Empirical Rule}
				\footnotesize
				$Z = (X-\mu)/\sigma \sim N(0,1)$. Empirical rule: $P(|X-\mu| \le 1\sigma) \approx 68\%$, $P(|X-\mu| \le 2\sigma) \approx 95\%$, $P(|X-\mu| \le 3\sigma) \approx 99.7\%$.
			\end{alertblock}
		\end{column}
	\end{columns}
\end{frame}

% SLIDE 04: Standardization
\begin{frame}{Z-Score Standardization}
\scriptsize
For $X\sim N(\mu,\sigma^2)$, the standardized variable
\[
Z=\frac{X-\mu}{\sigma}\sim N(0,1),\qquad
P(a<X<b)=P\!\left(\frac{a-\mu}{\sigma}<Z<\frac{b-\mu}{\sigma}\right).
\]
This transformation puts every Normal model on the same reference scale.
\end{frame}

% SLIDE 04B: MGF and CLT
\begin{frame}{MGF, Sum-Closure, and Central Limit Theorem}
	\vspace{-0.15cm}
	\begin{columns}[T]
		\begin{column}{0.48\textwidth}
			\begin{block}{MGF and Sum-Closure}
				\footnotesize
				$M_{X}(t) = e^{\mu t + \sigma^{2}t^{2}/2}$, with $M'(0)=\mu$ and $M''(0)=\mu^{2}+\sigma^{2}$. Sum of independent normals is normal: $N(\mu_{1}, \sigma_{1}^{2}) + N(\mu_{2}, \sigma_{2}^{2}) \sim N(\mu_{1}+\mu_{2}, \sigma_{1}^{2}+\sigma_{2}^{2})$.
			\end{block}
		\end{column}
		\begin{column}{0.48\textwidth}
			\begin{alertblock}{Central Limit Theorem}
				\footnotesize
				For i.i.d. with finite mean and variance, $Z_{n} = (\bar{X}-\mu)/(\sigma/\sqrt{n}) \xrightarrow{d} N(0,1)$. Universality: limit distribution is normal regardless of original distribution.
			\end{alertblock}
		\end{column}
	\end{columns}
\end{frame}

% SLIDE 04C: Empirical Rule
\begin{frame}{Empirical Rule: 68--95--99.7}
\scriptsize
For a Normal variable,
\[
P(|X-\mu|<\sigma)\approx0.68,\quad
P(|X-\mu|<2\sigma)\approx0.95,\quad
P(|X-\mu|<3\sigma)\approx0.997.
\]
These bands provide a fast diagnostic for unusual observations and process
variation before a formal tail calculation.
\end{frame}

% SLIDE 04D: Central Limit Theorem
\begin{frame}{Central Limit Theorem}
\scriptsize
If $X_1,\ldots,X_n$ are i.i.d. with finite mean $\mu$ and variance $\sigma^2$,
then
\[
\frac{\bar X-\mu}{\sigma/\sqrt n}\xrightarrow{d}N(0,1)
\qquad(n\to\infty).
\]
The limit explains why Normal approximations arise even when the original data
are not Normal.
\end{frame}

% SLIDE 05: Applications
\begin{frame}{Applications in Inference, Science, and Data Science}
	\vspace{-0.1cm}
	\begin{itemize}\itemsep=0.1cm
		\item \textbf{Statistical Inference:} $Z$-tests, confidence intervals, linear regression, and ANOVA assume Normality. \pause
		\item \textbf{Quality and Six Sigma:} $\bar{x}$ and $R$ control charts based on $\pm 3\sigma$ of the Normal. \pause
		\item \textbf{Finance:} Black-Scholes model for option pricing; geometric Brownian motions. \pause
		\item \textbf{Machine Learning:} Gaussian linear regression, generalized linear models, and likelihood function approximation.
	\end{itemize}
\end{frame}

% SLIDE 06: Python Lab Block 1
\begin{frame}[fragile]{Python Lab: PDF and Standardization}
	\vspace{-0.05cm}
	\scriptsize
	Validation of $\int f = 1$, Z-score standardization, and 68-95-99.7 empirical rule (\texttt{04.05\_normal\_distribution.py}):
	\vspace{0.02cm}
	{\lstset{basicstyle=\fontsize{4.5pt}{5.5pt}\selectfont\ttfamily}
	\lstinputlisting[firstline=15, lastline=33]{../../code/04_variables_aleatorias_continuas/04.05_normal_distribution.py}}
\end{frame}

% SLIDE 07: Python Lab Block 2
\begin{frame}[fragile]{Python Lab: Z-Test and Quantiles}
	\vspace{-0.05cm}
	\scriptsize
	Interval probabilities, quantiles, and one-sample Z-test:
	\vspace{0.02cm}
	{\lstset{basicstyle=\fontsize{4pt}{5pt}\selectfont\ttfamily}
	\lstinputlisting[firstline=40, lastline=63]{../../code/04_variables_aleatorias_continuas/04.05_normal_distribution.py}}
\end{frame}

% SLIDE 08: Python Lab Block 3
\begin{frame}[fragile]{Python Lab: CLT and Binomial-Normal Approximation}
	\vspace{-0.05cm}
	\scriptsize
	Monte Carlo verification of the CLT and application of the Binomial-Normal approximation with Yates correction:
	\vspace{0.02cm}
	{\lstset{basicstyle=\fontsize{4pt}{5pt}\selectfont\ttfamily}
	\lstinputlisting[firstline=73, lastline=100]{../../code/04_variables_aleatorias_continuas/04.05_normal_distribution.py}}
\end{frame}

% SLIDE 09: Python Lab Terminal Output
\begin{frame}[fragile]{Python Lab: Terminal Output}
	\vspace{-0.1cm}
	\begin{block}{}
		\fontsize{4pt}{5pt}\selectfont
		\begin{verbatim}
=== Block 1: Normal PDF & Standardization ===
Normal(100,15) PDF integral: 1.00000000
Z-score: F(x) actual = Phi(Z) (diff = 0)
Empirical rule: 68.27% | 95.45% | 99.73%

=== Block 2: Z-Score & Interval Probabilities ===
Heights N(170,100): P(X>185)=0.0668, P(160<=X<=180)=0.6827
Quantiles: q_0.025=-1.96, q_0.05=-1.64, q_0.5=0, q_0.95=1.64, q_0.975=1.96
Z-test: Z=1.5, p-value=0.1336, fail to reject H0

=== Block 3: CLT & Normal Approximation ===
Sum of 30 U(0,1) ~ N(15.0, 2.5): KS D=0.0014, p=0.71
Binomial(400, 0.25) approx:
  Exact: P(X>=120)=0.0133
  Normal w/ Yates: 0.0122 (error 0.0011)
  Required n for SE<0.01: 1875
		\end{verbatim}
	\end{block}
\end{frame}









% SLIDE 18: Closing and Next Steps
\begin{frame}{Closing Section and Modular Perspective}
	\vspace{-0.2cm}
	\begin{block}{Synthesis: The Gaussian Bell as Universal Distribution}
		\scriptsize
		The Normal distribution $N(\mu, \sigma^{2})$ is arguably the most important in all of statistics. Its bell shape, parameterized by only two moments, captures the essence of many natural phenomena. The Central Limit Theorem establishes its universality: the sum of i.i.d. random variables converges to the Normal regardless of the original distribution.
	\end{block}
	\vspace{0.05cm}
	\pause
	\begin{alertblock}{Key Lessons and Next Step}
		\begin{itemize}\itemsep=0.05cm
			\item \textbf{PDF:} $f(x) = \frac{1}{\sigma\sqrt{2\pi}} e^{-(x-\mu)^{2}/(2\sigma^{2})}$ models errors and fluctuations.
			\item \textbf{Standardization:} $Z = (X-\mu)/\sigma$ reduces to $N(0,1)$; probabilities are read from universal tables.
			\item \textbf{CLT:} $Z_{n} = (\bar{X}-\mu)/(\sigma/\sqrt{n}) \xrightarrow{d} N(0,1)$ explains the omnipresence of the Normal.
			\item \textbf{Next Section:} \textcolor{TecRojo}{Gamma, Beta, and Weibull Distributions}.
		\end{itemize}
	\end{alertblock}
\end{frame}

% SLIDE 20: Modular Perspective
\begin{frame}{Modular Perspective --- The Next Challenge}
\scriptsize
\begin{alertblock}{Next Session: Gamma, Beta, and Weibull}
The Normal family handles symmetric continuous variation. The next section adds
skewed waiting-time, proportion, and reliability models through Gamma, Beta, and
Weibull distributions.
\end{alertblock}
\begin{block}{Recommended Preparation}
Review standardization, the empirical rule, and the Central Limit Theorem.
\end{block}
\end{frame}

\end{document}
